% ============================================================
% §3 模型建立 — 数学骨架 (Stage 02/03 产物)
% 竞赛: 电工杯 B题 (2026)
% 阶段: Stage 02 问题深度解析 + Stage 03 模型选型
% ============================================================

\section{模型建立}

\subsection{全局符号说明}

为消除跨子问题的符号歧义，表\ref{tab:notation}给出全局符号约定。

\begin{table}[htbp]
\centering
\caption{全局符号说明表}
\label{tab:notation}
\small
\begin{tabular}{lll}
\toprule
\textbf{符号} & \textbf{含义} & \textbf{值/单位} \\
\midrule
$\mathcal{I},i,j$ & 小区索引集 $\{A,\dots,J\}$，$|\mathcal{I}|=10$ & -- \\
$\mathcal{K},k$ & 服务站规模 \{小,中,大\} & -- \\
$\mathcal{E}$ & 老人类型 \{自理,半失能,失能\} & -- \\
$\mathcal{S}$ & 服务项目 \{1,\dots,6\} & -- \\
\midrule
$\mathbf{S}_i^{(t)}$ & $i$小区第$t$年末三类老人数列向量 & 人 \\
$\mathbf{A}$ & 增生型状态转移矩阵 ($3\times 3$) & -- \\
$\alpha_{\text{S}\to\text{M}}$, $\alpha_{\text{M}\to\text{D}}$ & 自理$\to$半失能、半失能$\to$失能转移概率 & 0.045, 0.10 \\
$\mu$, $\beta$ & 死亡率、新增老年人占比 & 0.05, 0.07 \\
\midrule
$y_{i,k}$ & 在$i$建规模$k$服务站的0-1变量 & $\{0,1\}$ \\
$x_{i,j}$ & $j$小区是否由$i$站服务的0-1变量 & $\{0,1\}$ \\
$C_k^{\text{build}}$ & 规模$k$建设成本 & $\{18,32,45\}$ 万元 \\
$Q_k^{\max}$ & 规模$k$日最大服务人次 & $\{1000,2000,3000\}$ \\
$d_{i,j}$ & 小区间距离矩阵 & m \\
$B$ & 建设总预算约束 & 120 万元 \\
\midrule
$S_{i,j}$ & 综合满意度 $=0.2S1_{ij}+0.3S2_i+0.5S3_i$ & $[0.6, 1.0]$ \\
$S1_{i,j}$ & 距离满意度 (4段阶梯) & $\{0.60,0.75,0.90,1.00\}$ \\
$S2_i$ & 响应满意度 (5段阶梯) & $\{0.60,0.72,0.85,0.93,1.00\}$ \\
$S3_{i,s}$ & 价格满意度 (4段阶梯) & $\{0.60,0.75,0.90,1.00\}$ \\
\bottomrule
\end{tabular}
\end{table}

\subsection{Q1：老年人口演化的双模型互证体系}

为响应多模型交叉验证（Triangulation）的方法论原则，本文同时构造马尔可夫链与Leslie矩阵两套独立的人口演化模型，以数值结论的高度一致性作为预测可信度的数理担保。

\subsubsection{模型A：马尔可夫状态转移矩阵的物理构造}

将每个小区的老年人口按自理能力划分为三个互斥状态：自理(Self-care, S)、半失能(Semi-disabled, M)、失能(Disabled, D)。记第$t$年末小区$i$的老年人口状态向量为：

\begin{equation}
\mathbf{S}_i^{(t)} = \begin{bmatrix} s_{i,\text{S}}^{(t)} \\ s_{i,\text{M}}^{(t)} \\ s_{i,\text{D}}^{(t)} \end{bmatrix}, \quad i \in \mathcal{I}, \quad t = 0,1,\dots,5
\end{equation}

其中$t=0$对应于题目给出的基年截面数据。

题设动力学受三个独立机制支配：(1) 状态退化——自理$\xrightarrow{0.045}$半失能$\xrightarrow{0.10}$失能（吸收壁）；(2) 自然死亡——年均$\mu=0.05$概率退出；(3) 新生注入——每年新增$\beta=0.07$（当年总量）的60岁自理老人。编码为增生型转移矩阵$\mathbf{A}$：

\begin{equation}
\boxed{
\mathbf{A} =
\begin{bmatrix}
(1-\mu)(1-\alpha_{\text{S}\to\text{M}}) + \beta & \beta & \beta \\
(1-\mu)\alpha_{\text{S}\to\text{M}} & (1-\mu)(1-\alpha_{\text{M}\to\text{D}}) & 0 \\
0 & (1-\mu)\alpha_{\text{M}\to\text{D}} & (1-\mu)
\end{bmatrix}
}
\end{equation}

代入$\mu=0.05$, $\alpha_{\text{S}\to\text{M}}=0.045$, $\alpha_{\text{M}\to\text{D}}=0.10$, $\beta=0.07$得数值矩阵，列和$1.02>1$反映年均净增长约2\%——"增生型马尔可夫链"区别于标准吸收链的核心特征。

\subsubsection{模型B：Leslie矩阵互证}

将三类老人视为三个"生理年龄组"，构造$3\times 3$ Leslie矩阵$\mathbf{L}$\cite{leslie1945}：$f_1=(1-\mu)(1-\alpha_{\text{S}\to\text{M}})+\beta$, $f_2=f_3=\beta$, $p_{1\to2}=(1-\mu)\alpha_{\text{S}\to\text{M}}$, $p_{2\to3}=(1-\mu)\alpha_{\text{M}\to\text{D}}$。$\mathbf{L}\equiv\mathbf{A}^\top$，主特征值$\lambda_1=1.0198>1$揭示年均净增长率约2.0\%，与Q1数值结果精确吻合——纯Markov递推无法直接提供此结构稳定性信息（数值对比见附录表\ref{tab:leslie_check}）。

\subsubsection{递推预测算法}

\begin{algorithm}[htbp]
\caption{老年人口5年递推预测}
\label{alg:markov}
\begin{algorithmic}[1]
\Require 基年状态向量 $\mathbf{S}_i^{(0)}$ (取自题目数据), 转移矩阵 $\mathbf{A}$, 预测年限 $T=5$
\Ensure 各年末状态向量 $\{\mathbf{S}_i^{(t)}\}_{t=1}^{T}$
\For{$t = 0$ \textbf{to} $T-1$}
    \For{$i \in \mathcal{I}$}
        \State $\mathbf{S}_i^{(t+1)} \gets \mathbf{A} \cdot \mathbf{S}_i^{(t)}$
        \State 各分量取整: $s_{i,e}^{(t+1)} \gets \lfloor s_{i,e}^{(t+1)} + 0.5 \rfloor$
    \EndFor
\EndFor
\Return $\{\mathbf{S}_i^{(t)}\}_{t=1}^{T}$
\end{algorithmic}
\end{algorithm}

\subsubsection{服务需求量的分层预测 (Q1.2--Q1.3)}

第$t$年末小区$i$类型$e$老人对服务$s$的\textbf{理论月需求次数}为：

\begin{equation}
\hat{D}_{i,e,s}^{(t)} = s_{i,e}^{(t)} \cdot D_{e,s}, \quad \forall i\in\mathcal{I}, e\in\mathcal{E}, s\in\mathcal{S}
\end{equation}

其中$D_{e,s}$为题目给出的每位老人月均服务需求次数。

Q1.3引入消费约束后，定义类型$e$老人的月均全量服务消费额：

\begin{equation}
C_e^{\text{full}} = \sum_{s\in\mathcal{S}} D_{e,s} \cdot p_s
\end{equation}

若$C_e^{\text{full}} \le w_i \cdot M_e$，则消费约束不激活，实际需求等于理论需求；否则触发题目规定的等比例削减规则：

\begin{equation}
D_{i,e,s}^{\text{actual}} = s_{i,e}^{(t)} \cdot D_{e,s} \cdot \min\left(1, \frac{w_i \cdot M_e}{C_e^{\text{full}}}\right), \quad \text{取整}
\end{equation}

由§2.3的EDA预诊断知，失能老人在全部10个小区均触发削减（缺口23.9\%--49.1\%），半失能老人在2个低收小区触发，自理老人均不触发——此信息将直接写入Q1.3的数值计算结果。

\subsection{Q2：嵌入式选址-规模联合优化的MILP模型}

\subsubsection{问题定位于MCLP-CFLP谱系交汇处}

Q2是经典\textbf{最大覆盖选址（MCLP, Church\&ReVelle, 1974\cite{church1974}）}与\textbf{有容量设施选址（CFLP\cite{daskin1995}）}的复合变体：1000m服务半径对应MCLP覆盖逻辑，三种规模的硬容量约束对应CFLP多容量等级，满意度$S$的分段阶梯函数带来了额外的非线性。本文将所建模型命名为\textbf{MILP-MCLP-ES}，通过Big-M线性化将非凸满意度规则转化为混合整数线性不等式组，以分支定界求全局最优。淘汰方案：GA/PSO（非凸可行域中无最优性保证）；AHP/TOPSIS（用评价替代优化，丧失空间-容量-满意度耦合信息）\cite{wang2024,sima2024}。

\subsubsection{目标函数}

采用线性加权法将双目标聚合为单目标：

\begin{equation}
\max \; \lambda \cdot \frac{\sum_{i\in\mathcal{I}}\sum_{j\in\mathcal{I}}\sum_{e\in\mathcal{E}} u_{i,j,e} \cdot s_{j,e}^{(5)}}{\sum_{j\in\mathcal{I}}\sum_{e\in\mathcal{E}} s_{j,e}^{(5)}} + (1-\lambda) \cdot \frac{\sum_{i\in\mathcal{I}}\sum_{j\in\mathcal{I}} S_{i,j} \cdot x_{i,j}}{\sum_{i\in\mathcal{I}}\sum_{j\in\mathcal{I}} x_{i,j}}
\end{equation}

其中$\lambda \in [0,1]$为帕累托权重。第一项为服务覆盖率（享受服务的老年人数占比）；第二项为平均综合满意度。后期通过变动$\lambda \in \{0.3, 0.5, 0.7\}$绘制帕累托前沿。

\subsubsection{约束条件}

\textbf{(C1) 建设预算约束}：
\begin{equation}
\sum_{i\in\mathcal{I}}\sum_{k\in\mathcal{K}} C_k^{\text{build}} \cdot y_{i,k} \le B = 120 \; \text{(万元)}
\end{equation}

\textbf{(C2) 每小区至多建一个服务站}：
\begin{equation}
\sum_{k\in\mathcal{K}} y_{i,k} \le 1, \quad \forall i \in \mathcal{I}
\end{equation}

\textbf{(C3) 服务半径约束}：只有距离$\le 1000$m的小区对才可能产生服务关系。
\begin{equation}
x_{i,j} \le \sum_{k\in\mathcal{K}} y_{i,k}, \quad \forall i,j \in \mathcal{I}
\end{equation}
\begin{equation}
x_{i,j} = 0, \quad \forall i,j: d_{i,j} > 1000 \text{ m}
\end{equation}

\textbf{(C4) 唯一分配约束}：每位老人由恰好一个服务站服务（满意度最大化$\equiv$距离最小化，因Q2阶段定价固定$S_3\equiv1.0$）。
\begin{equation}
\sum_{i\in\mathcal{I}} x_{i,j} = 1, \quad \forall j \in \mathcal{I}
\end{equation}

\textbf{(C5) 服务容量约束}：
\begin{equation}
\sum_{j\in\mathcal{I}}\sum_{e\in\mathcal{E}} \sum_{s\in\mathcal{S}} D_{i,j,e,s}^{\text{actual}} \cdot x_{i,j} \le \sum_{k\in\mathcal{K}} Q_k^{\max} \cdot y_{i,k} \cdot 30, \quad \forall i \in \mathcal{I}
\end{equation}
其中30为月均天数（将日容量转为月容量），$D_{i,j,e,s}^{\text{actual}}$为Q1.3求得的实际需求。

\textbf{(C6) 站点存在性约束}：
\begin{equation}
x_{i,i} \ge y_{i,1} + y_{i,2} + y_{i,3}, \quad \forall i \in \mathcal{I}
\end{equation}
若在小区$i$建站，该小区老人必须由本站服务（距离=0，满意度最高）。

\subsubsection{满意度规则的精确线性化：Big-M + SOS2方法}

题目定义的满意度函数$S = 0.2S_1 + 0.3S_2 + 0.5S_3$是三个分段常数函数的加权和。以下逐一给出其MILP线性化编码。

\paragraph{距离满意度 $S_1$ 的线性化}

$S_1(d)$为4档阶梯函数，断点$\{0,300,500,650,1000\}$m，分值$\{1.00,0.90,0.75,0.60\}$。引入$\delta_{i,j,\ell}^{S1}\in\{0,1\}$ ($\ell=1,\dots,4$)：

\begin{equation}
\begin{aligned}
\sum_{\ell=1}^{4} \delta_{i,j,\ell}^{S1} &= x_{i,j}, \quad d_{i,j} \le 300\delta_1 + 500\delta_2 + 650\delta_3 + M_d\delta_4 \\
d_{i,j} &\ge 0\cdot\delta_1 + 300\delta_2 + 500\delta_3 + 650\delta_4 \\
S1_{i,j} &= 1.00\delta_1 + 0.90\delta_2 + 0.75\delta_3 + 0.60\delta_4
\end{aligned}
\end{equation}

\paragraph{响应满意度 $S_2$ 的线性化}

$S_2(\rho)$为5档阶梯函数，断点$\{0,0.60,0.75,0.85,0.95,1.00\}$，分值$\{1.00,0.93,0.85,0.72,0.60\}$。利用率$\rho_i$含分式，引入$\delta_{i,\ell}^{S2}\in\{0,1\}$与McCormick包络辅助变量$w_{ij}=S2_i\cdot x_{ij}$消去非线性乘积项\cite{mccormick1976}：

\begin{equation}
\begin{aligned}
\sum_{\ell=1}^{5} \delta_{i,\ell}^{S2} &= \sum_{k} y_{i,k}, \quad S2_i = 1.00\delta_1 + 0.93\delta_2 + 0.85\delta_3 + 0.72\delta_4 + 0.60\delta_5 \\[4pt]
w_{ij} &\le x_{ij},\; w_{ij} \le S2_i,\; w_{ij} \ge S2_i - (1-x_{ij}),\; 0 \le w_{ij} \le 1
\end{aligned}
\end{equation}

\paragraph{价格满意度 $S_3$ 的线性化 (Q3用)}

$S_3$断点以基准价$p_s$为参照：$\{1.00p_s,1.10p_s,1.20p_s\}$，分值$\{1.00,0.90,0.75,0.60\}$。在Q2定价固定的情况下$S_3\equiv1.00$；Q3引入$\delta_{i,s,\ell}^{S3}$并附加词典序目标（优先$S_3$最大化、次优$\varepsilon=10^{-4}$营收激励），详见§5.3。

\paragraph{综合满意度}

\begin{equation}
S_{i,j} = 0.2 \cdot S1_{i,j} + 0.3 \cdot S2_i + 0.5 \cdot \frac{1}{|\mathcal{S}|}\sum_{s\in\mathcal{S}} S3_{i,s}, \quad \forall i,j \in \mathcal{I}: x_{i,j}=1
\end{equation}

\subsubsection{线性化变量规模}

上述Big-M方案的二元变量总数$|\mathcal{I}|^2\cdot4 + |\mathcal{I}|\cdot5 + |\mathcal{I}|\cdot|\mathcal{S}|\cdot4 = 690$，加上$y_{i,k}$(30个)和$x_{i,j}$(100个)，总计约820个二元变量、约束约2000条，CBC分支定界可在秒级收敛\cite{wolsey1998}。

\subsection{Q3：服务定价与政府补贴优化模型}

Q3在Q2最优选址$\{y_{i,k}^*\}$固化的前提下，允许各站自主定价$p'_{i,s}$，目标最大化满意度，受利润率$\le 8\%$监管约束\cite{feng2024,zhao2024}。定价通过$S_3$影响满意度（$x_{i,j}^*$固定）。

\subsubsection{利润率与补贴约束}

年利润$\Pi_i$含运营收支、政府补贴$\sigma_s^{\text{sub}}=2$元/人次（紧急救助除外）、20年折旧：

\begin{equation}
\Pi_i = 360 \sum_{j,e,s} D_{i,j,e,s}^{\text{actual}} (p'_{i,s} - c_s + \sigma_s^{\text{sub}}) x_{i,j}^* - 360 \sum_k O_k y_{i,k}^* - \frac{\sum_k C_k^{\text{build}} y_{i,k}^*}{20}
\end{equation}

\begin{equation}
\frac{\Pi_i}{360 \sum_k O_k y_{i,k}^* + \sum_k C_k^{\text{build}} y_{i,k}^* / 20} \le 0.08, \quad \forall i
\end{equation}

补贴上限约束（非紧急救助）：$\sum_{j,e,s} D_{i,j,e,s}^{\text{actual}} \cdot 2 \cdot x_{i,j}^* \le B_k^{\text{sub}}, \; B_k^{\text{sub}} \in \{1000, 1800, 2600\}$元/日。

\subsection{模型族一致性}

Q1$\to$Q2$\to$Q3串行依赖链：Q1输出人口与需求$\to$Q2选址$\to$Q3定价。Q2与Q3共享Big-M满意度线性化框架，Q4通过参数扰动复用Q2/Q3求解器。全文方法内聚于统一的MILP编译框架。
